% Part: second-order-logic
% Chapter: sol-and-sets
% Section: introduction

\documentclass[../../../include/open-logic-section]{subfiles}

\begin{document}

\olfileid{sol}{set}{int}

\olsection{Pambuka}

Amarga logika orde kapindho bisa nguwantifikasi
subhimpunan domain uga fungsi, kita bisa ngarepake yen
saperangan teori himpunan bisa diandharake ing logika
orde kapindho. Kang diarani ``diandharake'' ing kene
yaiku sipat lan pratelan teori himpunan bisa dinyatakake
ing logika orde kapindho tanpa kosakata nonlogis kang
mligi kanggo himpunan (umpamane !!{predicate}
kang nyatakake dadi anggota ing teori himpunan). Umpamane, kita bisa
netepake gabungan lan irisan himpunan sarta relasi
subhimpunan, uga mbandhingake ukuran himpunan lan
nyatakake asil kaya Teorema Cantor.

\end{document}
